\documentclass[11pt,aspectratio=169]{beamer}
\usetheme{Boadilla}
\usefonttheme{professionalfonts}
\setbeamertemplate{navigation symbols}{}

\usepackage[T1]{fontenc}
\usepackage{lmodern}
\usepackage{amsmath,amssymb,mathtools}
\usepackage{booktabs}
\usepackage{tabularx}
\usepackage{array}
\usepackage{graphicx}
\usepackage{xcolor}

\graphicspath{{figs/}}

% Show section in a small headline so the audience tracks the arc
\setbeamertemplate{headline}{%
  \hbox{}\vskip2pt\hfill\textcolor{gray!70}{\scriptsize\insertsectionhead}\hspace{1em}\vskip-6pt%
}
\setbeamertemplate{section in toc}[sections numbered]
\setbeamertemplate{frametitle}[default][left]
\setbeamerfont{framesubtitle}{size=\small}

\newcommand{\figplaceholder}[1]{\begin{center}\fbox{\begin{minipage}[c][0.32\textheight][c]{0.82\textwidth}\centering\large #1\end{minipage}}\end{center}}
\newcommand{\shortfigplaceholder}[1]{\begin{center}\fbox{\begin{minipage}[c][0.22\textheight][c]{0.82\textwidth}\centering\large #1\end{minipage}}\end{center}}
\newcommand{\compactfigplaceholder}[1]{\begin{center}\fbox{\begin{minipage}[c][0.10\textheight][c]{0.82\textwidth}\centering\normalsize #1\end{minipage}}\end{center}}
\newcommand{\R}{\mathcal{R}}
\newcommand{\E}{\mathcal{E}}
\newcommand{\dd}{\,\mathrm{d}}
\newcommand{\good}[1]{\textcolor{blue!70!black}{#1}}
\newcommand{\bad}[1]{\textcolor{red!70!black}{#1}}
\newcommand{\app}[1]{\hyperlink{#1}{\beamergotobutton{details}}}
\newcommand{\back}[1]{\hyperlink{#1}{\beamerreturnbutton{back}}}

\title[Directional Congestion]{Directional Congestion}
\subtitle{Measurement, Spatial Equilibrium, and Tidal-Lane Reallocation}
\author[Duan, Feng, Gu]{Yige Duan (SJTU) $\quad$ Xiaoruo Feng (PKU) $\quad$ Yizhen Gu (PKU)}
\date{\today}

\begin{document}

\begin{frame}
\titlepage
\end{frame}

\begin{frame}{Roadmap}
\footnotesize\tableofcontents
\end{frame}

% =====================================================================
\section{Motivation}
% =====================================================================
\begin{frame}{Motivation}
\begin{itemize}
  \item \textbf{Peak-direction imbalance.}
  \begin{itemize}\footnotesize
    \item The two directions of the \emph{same} physical road carry very different traffic at the same time.
    \item Where workers live vs.\ where jobs sit pins down which direction floods.
    \item Same lanes, symmetric pavement, yet \bad{directional} congestion.
  \end{itemize}
  \item \textbf{The daily round trip.}
  \begin{itemize}\footnotesize
    \item Morning and evening reverse the same flows.
    \item Value-of-time costs should count both legs.
  \end{itemize}
  \item \textbf{A cheap policy margin.}
  \begin{itemize}\footnotesize
    \item \good{Tidal (reversible) lanes} shift one lane toward the busy direction at \good{zero net lane-km}.
    \item They can still move route choice \emph{and} residence--workplace sorting.
  \end{itemize}
\end{itemize}
\vspace{0.2em}
\compactfigplaceholder{Figure 1 (to insert): AM directional speed / traffic-per-lane asymmetry on the study network}
\end{frame}

% =====================================================================
\section{This Paper}
% =====================================================================
\begin{frame}{This Paper}
\begin{itemize}
  \item \textbf{Model directional congestion in a spatial GE.}
  \begin{itemize}\footnotesize
    \item Extend the Allen--Arkolakis inner-city model.
    \item Keep each road \good{physically symmetric} ($t_{kl}^{\mathrm{ff}}=t_{lk}^{\mathrm{ff}}$) but let \good{traffic-per-lane} differ by direction, so realized $t_{kl}\neq t_{lk}$ and $\tau_{ij}\neq\tau_{ji}$ endogenously.
  \end{itemize}
  \item \textbf{Evaluate tidal lanes as a capacity-neutral instrument.}
  \begin{itemize}\footnotesize
    \item Top-100 reciprocal pairs, $+1$ lane to the slow direction, $-1$ to the fast direction, total capacity fixed.
  \end{itemize}
  \item \textbf{Decompose the effect into three transparent channels.}
  \begin{itemize}\footnotesize
    \item \good{ST engineering gain} (fixed sorting) $\to$ \good{network GE wedge} $\to$ \bad{residence--workplace sorting}.
  \end{itemize}
  \item \textbf{Headline.}
  \begin{itemize}\footnotesize
    \item ST \good{saves time and raises welfare}; long-run \bad{sorting capitalizes the gain into market access}.
    \item Commute time nets out near zero; aggregate welfare turns slightly negative, with the loss running through $\Pi$ rather than amenities.
  \end{itemize}
\end{itemize}
\end{frame}

% =====================================================================
\section{Literature}
% =====================================================================
\begin{frame}{Literature}
\begin{itemize}
  \item \textbf{Spatial equilibrium and agglomeration.}
  \begin{itemize}\footnotesize
    \item Allen and Arkolakis (RES, 2022); Ahlfeldt et al.\ (2015).
  \end{itemize}
  \item \textbf{Transportation policy and congestion.}
  \begin{itemize}\footnotesize
    \item Barwick et al.\ (AER, 2024); Yang et al.\ (AEJ:\ Policy, 2020).
  \end{itemize}
  \item \textbf{Calibration and contribution.}
  \begin{itemize}\footnotesize
    \item Beijing evidence gives a congestion elasticity $\lambda\approx1.1$ (edge time rises with traffic per lane).
    \item We add a \underline{tidal-lane margin} and a full ST / GE / sorting welfare decomposition under endogenous \emph{directional} congestion.
  \end{itemize}
\end{itemize}
\end{frame}

% =====================================================================
\section{Data}
% =====================================================================
\begin{frame}[label=data]{Data}
\small
\begin{columns}[T,onlytextwidth]
\column{0.58\textwidth}
\begin{itemize}
  \item \textbf{Travel cost (Map API).}
  \begin{itemize}\footnotesize
    \item Directional travel time and speed, crawled by direction and time of day; the core directional-cost input.
    \item Collected at several scales: national township-pair OD, city grids (Shenzhen 1\,km), and Beijing road centerlines.
  \end{itemize}
  \item \textbf{Congestion index.}
  \begin{itemize}\footnotesize
    \item City-level congestion ranking (Baidu, 2025), used to relate congestion to asymmetry.
  \end{itemize}
  \item \textbf{Roads and capacity.}
  \begin{itemize}\footnotesize
    \item Road geometry, length, lanes, and functional class.
  \end{itemize}
  \item \textbf{Commuting.}
  \begin{itemize}\footnotesize
    \item Residence, workplace, OD flows, and directional travel demand.
  \end{itemize}
  \item \textbf{Estimation sample.}
  \begin{itemize}\footnotesize
    \item $\sim$1{,}100 grid nodes, $\sim$2{,}600 directed edges, $\sim$88{,}000 commuters, AM peak.
  \end{itemize}
\end{itemize}
\column{0.40\textwidth}
\compactfigplaceholder{Figure (to insert): study area and main road network}
\end{columns}
\vspace{0.2em}
\centering\app{appdata}\; sources, crawling, and construction
\end{frame}

% =====================================================================
\section{Pattern: Directional Asymmetry}
% =====================================================================
\begin{frame}{Pattern: A Ladder of Evidence}
\framesubtitle{Directional congestion is a \good{general urban feature}, not a Shenzhen tidal-lane special case.}
\footnotesize
\begin{tabularx}{\textwidth}{@{}p{0.30\textwidth}X@{}}
\toprule
Scale & What it shows \\
\midrule
1. National city OD & Asymmetry is \emph{widespread} across Chinese cities \\
2. Congested-city ranking & Congested cities are \emph{more} directional; AM/PM can reverse \\
3. \good{Shenzhen grid OD (bridge)} & Within one city, \good{real commute flow} is exposed to stronger asymmetry \\
4. Beijing road / grid & Asymmetry exists at the \emph{road-segment} level, not just OD aggregation \\
5. Beijing demand direction & Commuting pressure is itself directional, so symmetric cost mis-measures exposure \\
\midrule
($\ast$) Shenzhen corridor / tidal lane & A real policy already manages \emph{directional capacity} (illustration, not framing) \\
\bottomrule
\end{tabularx}
\vspace{0.4em}
\begin{itemize}\footnotesize
  \item Metric throughout: $\min(t_{ij},t_{ji})/\max(t_{ij},t_{ji})$, where a lower ratio means more asymmetric.
\end{itemize}
\end{frame}

\begin{frame}{Level 1. National Cities: Asymmetry Is Widespread}
\framesubtitle{Township-pair AM travel time, top-50 cities (2021).}
\begin{columns}[T,onlytextwidth]
\column{0.50\textwidth}
\scriptsize
\begin{tabular}{lccc}
\toprule
City & med.\ ratio & $\leq0.90$ & AM$>$night \\
\midrule
Beijing & \good{0.837} & 71.6\% & 77.3\% \\
Chengdu & 0.853 & 75.8\% & 63.4\% \\
Chongqing & 0.856 & 70.0\% & 66.6\% \\
Shenzhen & 0.864 & 63.3\% & 72.1\% \\
Xi'an & 0.873 & 65.6\% & 62.4\% \\
Guangzhou & 0.879 & 66.1\% & 65.8\% \\
Shanghai & 0.880 & 65.1\% & 79.1\% \\
\bottomrule
\end{tabular}
\vspace{0.3em}
\begin{itemize}\footnotesize
  \item Pooled weighted median \good{$0.903$}; \good{$52.2\%$} of weighted OD pairs $\leq0.90$.
  \item A typical two-way OD differs $\sim$10\% by direction, and AM is more asymmetric than night.
\end{itemize}
\column{0.49\textwidth}
\centering\scriptsize AM directional asymmetry by city\\
\includegraphics[width=\linewidth,height=0.6\textheight,keepaspectratio]{pat_national_city.png}
\end{columns}
\end{frame}

\begin{frame}{Level 2. Congested Cities Are More Directional}
\framesubtitle{Baidu 2025 congestion ranking vs.\ asymmetry (99 cities).}
\begin{center}\scriptsize
\begin{tabular}{lcccc}
\toprule
City & congestion rank & AM med. & PM med. & AM/PM reversal \\
\midrule
Beijing & 1 & \good{0.805} & 0.866 & 37.3\% \\
Chongqing & 2 & 0.790 & 0.872 & 46.1\% \\
Wuhan & 3 & 0.866 & 0.916 & 42.2\% \\
Guangzhou & 4 & 0.854 & 0.879 & 34.1\% \\
Shanghai & 5 & 0.862 & 0.894 & \good{66.7\%} \\
Chengdu & 9 & 0.767 & 0.889 & 47.2\% \\
Xi'an & 12 & 0.768 & 0.861 & 42.8\% \\
\bottomrule
\end{tabular}
\end{center}
\vspace{0.2em}
\begin{itemize}\footnotesize
  \item High-congestion cities are not just slow; they are slow \bad{with a direction}.
  \item A large share of OD pairs \emph{reverse} their slow direction between AM and PM.
  \item Directional congestion is a peak-timing phenomenon, exactly the tidal margin.
\end{itemize}
\end{frame}

\begin{frame}{Level 3. Shenzhen Grid: Real Flow Sees Stronger Asymmetry}
\framesubtitle{Shenzhen 1\,km fishnet two-way travel-time matrix; the cleanest within-city evidence.}
\begin{columns}[T,onlytextwidth]
\column{0.52\textwidth}
\scriptsize
\begin{tabular}{lccc}
\toprule
Weighting & med.\ ratio & $\leq0.90$ & $\leq0.83$ \\
\midrule
All grid pairs (1.46M) & 0.879 & 57.5\% & 35.1\% \\
\good{Observed commute flow} & \good{0.819} & \good{71.6\%} & 53.0\% \\
Car-asset weighted & 0.821 & 71.2\% & 52.5\% \\
\bottomrule
\end{tabular}
\vspace{0.4em}
\begin{itemize}\footnotesize
  \item \good{Flow-weighting makes asymmetry stronger}, not weaker: it overlaps with real commute \bad{exposure}.
  \item Heterogeneous within the city: workplace in econ-zone 1 has median $0.805$ ($73.7\%\leq0.90$).
  \item Caveat: grid500 flow-aggregate path (coverage $85.4\%$), not the raw 1\,km OD crosswalk.
\end{itemize}
\column{0.46\textwidth}
\centering\scriptsize All-pair vs.\ flow-weighted\\
\includegraphics[width=\linewidth,height=0.62\textheight,keepaspectratio]{pat_sz_floww.png}
\end{columns}
\end{frame}

\begin{frame}{Level 4. Beijing Road / Grid: Not an OD Artifact}
\begin{columns}[T,onlytextwidth]
\column{0.50\textwidth}
\begin{itemize}\small
  \item \textbf{Centerline segments (AM, 3{,}670 two-way).}
  \begin{itemize}\scriptsize
    \item Median ratio $\approx0.884$, $54.7\%\leq0.90$.
  \end{itemize}
\end{itemize}
\scriptsize
\begin{center}
\begin{tabular}{lcc}
\toprule
strict rule & AM segments & share \\
\midrule
unsym1 & 405 & 11.0\% \\
unsym2 & 179 & 4.9\% \\
unsym3 & 98 & 2.7\% \\
\bottomrule
\end{tabular}
\end{center}
\vspace{0.2em}
\begin{itemize}\small
  \item \textbf{Adjacent-grid (AM, 1{,}894 two-way).}
  \begin{itemize}\scriptsize
    \item Median ratio \good{$0.794$}, $69.7\%<0.90$.
  \end{itemize}
  \item Asymmetry shows up in \emph{local segments and neighboring cells}, a network mechanism rather than aggregation.
  \item \textbf{Quantification sample:} structural counterfactuals use the Beijing square grid (gate5 LCC; $\S$\ref{sec:quant}).
\end{itemize}
\column{0.49\textwidth}
\centering\scriptsize Directionally asymmetric AM segments (Beijing)\\
\includegraphics[width=\linewidth,height=0.66\textheight,keepaspectratio]{pat_bj_centerline.png}
\end{columns}
\end{frame}

\begin{frame}{Level 5. Directional Demand and a Real Tidal Lane}
\begin{columns}[T,onlytextwidth]
\column{0.52\textwidth}
\begin{itemize}\small
  \item \textbf{Demand is directional (Beijing).}
  \begin{itemize}\footnotesize
    \item \good{64.8\%} of car$+$bus weighted demand moves \emph{toward} the nearer job cluster; mean distance-to-center falls \good{$4.68$ km}.
    \item Some inbound sectors: inbound/reverse demand ratio $2$ to $6$.
    \item If demand is directional and congestion is direction-specific, \bad{a symmetric commuting cost mis-allocates exposure}.
  \end{itemize}
  \item \textbf{Policy already does this (Caitian Rd, Shenzhen).}
  \begin{itemize}\footnotesize
    \item Post-only: policy/anti-direction speed ratio $0.789\to\good{0.982}$; asymmetry DDD $>0$ vs.\ four control sets.
    \item Reverse-direction loss is more consistent than the treated gain, so it illustrates directional capacity management, not the main framing.
  \end{itemize}
\end{itemize}
\column{0.46\textwidth}
\centering\scriptsize Beijing directional commuting demand\\
\includegraphics[width=\linewidth,height=0.42\textheight,keepaspectratio]{pat_bj_demand.png}\\[2pt]
\scriptsize Caitian Rd event study (DiD)\\
\includegraphics[width=\linewidth,height=0.30\textheight,keepaspectratio]{pat_caitian.png}
\end{columns}
\end{frame}

% =====================================================================
\section{Model}
% =====================================================================
% --- Original model section (preserved verbatim from tidalraw.tex) ---

\begin{frame}{Model: Environment}
\begin{itemize}
  \item Locations $i,j\in N$ host $\bar{L}$ mass of residents and jobs
  \vspace{0.5em}
  \item A physical road is represented by two directed links, $(k,l)$ and $(l,k)$
  \begin{itemize}
    \item Symmetric \textbf{free-flow} time $t_{kl}^{ff}=t_{lk}^{ff}$, but directed lanes $n_{kl}$ and $n_{lk}$ can differ
    \item A tidal-lane policy shifts one lane across directions while holding $n_{kl}+n_{lk}$ fixed
  \end{itemize}
  \vspace{0.5em}
  \item The realized edge time $t_{kl}$ is \textcolor{blue}{endogenous} and \textcolor{red}{directed}
  \begin{itemize}
    \item $t_{kl}$ rises with traffic per lane $\Xi_{kl}/n_{kl}$
    \item $t_{kl} \neq t_{lk}$ when traffic differs by direction
  \end{itemize}
\end{itemize}
\end{frame}

\begin{frame}{Individual Location Choice}
\begin{itemize}
  \item An individual $\omega$ living in $i$, working in $j$, and taking route $r\in\R_{ij}$ obtains
  \vspace{-0.5em}\begin{footnotesize}\begin{equation}
    V_{ij,r}(\omega) = \left( \frac{u_i w_j}{C_{ij,r}} \right)\varepsilon_{ij,r}(\omega),
    \qquad
    C_{ij,r} = \prod_{\ell=1}^{K_r} t_{r_{\ell-1},r_\ell}
  \vspace{-1.25em}\end{equation}\end{footnotesize}
  \begin{itemize}
    \item $u_i$ is amenity; $w_j$ is wage; $\epsilon$ is i.i.d. Frechet
    \item $C_{ij,r}$ is the route cost (product of edge times); route $r=[r_0,\dots,r_{K_r}]$, $r_0=i$, $r_{K_r}=j$
  \end{itemize}
  \vspace{0.5em}
  \item Choice probabilities are
  \vspace{-0.5em}\begin{equation}\label{eq:pi}
    \pi_{ij,r} = \frac{\left(u_i w_j/C_{ij,r}\right)^\theta}{\sum_{a}\sum_b\sum_{s\in\R_{ab}}\left(u_a w_b/C_{ab,s}\right)^\theta}
  \end{equation}
\end{itemize}
\end{frame}

\begin{frame}{Bilateral Commuting Costs and Flows}
\begin{itemize}
  \item The bilateral commuting-cost, aggregated across all routes (Frechet expectation)
  \vspace{-0.5em}\begin{equation}
    \tau_{ij} = \left[ \sum_{r\in\R_{ij}} C_{ij,r}^{-\theta} \right]^{-1/\theta}
  \end{equation}
  \item Thus, residence-workplace choice shares satisfy
  \vspace{-0.25em}\begin{equation}\label{eq:gravity}
    L_{ij} = \sum_{r\in\R_{ij}}L_{ij,r} = \tau_{ij}^{-\theta} \left(\frac{L_i^R}{\Pi_i}\right) \left(\frac{L_j^F}{P_j}\right), \,\, L_i^R=\sum_jL_{ij}, \,\, L_j^F=\sum_iL_{ij}
  \vspace{-1em}\end{equation}
  \begin{itemize}
    \item Directional congestion implies $\tau_{ij}\neq\tau_{ji}$ in general
    \item $\Pi_i$ and $P_j$ are origin and destination market-access terms
  \end{itemize}
\end{itemize}
\end{frame}

\begin{frame}{Amenity, Productivity, and Location Choices}
\begin{itemize}
  \item Agglomeration of local fundamentals
  \vspace{-0.5em}\begin{equation}\label{eq:agglom}
    w_j = A_j = \bar A_j\left(L_j^F\right)^\alpha, \qquad u_i=\bar u_i\left(L_i^R\right)^\beta
  \end{equation}
  \item Resident ($R$) and workplace ($F$) distribution across space (level-form shares; v4 eqs.\ invert\_R--invert\_F)
  \begin{align}
  (l_i^R)^{1-\beta\theta} & = \chi \sum_j\tau_{ij}^{-\theta}\,\bar u_i^{\theta}\,\bar A_j^{\theta}\,(l_j^F)^{\alpha\theta} \label{eq:clear1} \\
  (l_j^F)^{1-\alpha\theta} & = \chi \sum_i\tau_{ij}^{-\theta}\,\bar u_i^{\theta}\,\bar A_j^{\theta}\,(l_i^R)^{\beta\theta} \label{eq:clear2}
  \end{align}
  \begin{itemize}\footnotesize
    \item Mainline: $\theta=6.83$, $\alpha=+0.05$, $\beta=-0.10$ $\Rightarrow$ exponents $\approx 1.66$ (residence) and $\approx 0.66$ (workplace).
  \end{itemize}
\end{itemize}
\end{frame}

\begin{frame}{From Edge Times to Bilateral Commuting Cost}
\begin{itemize}
  \item Define the directed adjacency matrix $A=[a_{kl}]$ by
  \vspace{-0.5em}\begin{equation}
    a_{kl} = t_{kl}^{-\theta}
  \end{equation}
  \item Summing over all feasible route lengths gives
  \vspace{-0.5em}\begin{equation}
    B = \sum_{K=0}^{\infty}A^K = (I-A)^{-1}
  \end{equation}
  \item The bilateral commuting-cost aggregator can be written as
  \vspace{-0.5em}\begin{equation}\label{eq:inverse}
    \tau_{ij} = B_{ij}^{-1/\theta}
  \end{equation}
\end{itemize}
\end{frame}

\begin{frame}{Feedback between Location Choice, Traffic, and Edge Time}
\begin{itemize}
  \item For a directed trip $i\to j$, define $\kappa^{kl}_{ij}$ as expected use of link $(k,l)$
  \vspace{-1em}\begin{equation}
    \kappa^{kl}_{ij} = \left( \frac{\tau_{ij}} {\tau_{ik}\,t_{kl}\,\tau_{lj}} \right)^\theta
  \vspace{-0.5em}\end{equation}
  \item The directional traffic is
  \begin{equation}\label{eq:Xi_c}
  \Xi_{kl} = \sum_i\sum_j \kappa^{kl}_{ij}L_{ij}
  \end{equation}
  \item Traffic, in turn, raises the directed edge time (\good{constant-elasticity / power form})
  \vspace{-0.5em}\begin{equation}\label{eq:c_Xi}
    t_{kl} = t_{kl}^{ff} \left(\frac{\Xi_{kl}}{n_{kl}}\right)^{\lambda},\qquad \lambda\approx 1.1
  \vspace{-0.5em}\end{equation}
    This yields asymmetric edge times $t_{kl}\neq t_{lk}$ even when $t_{kl}^{ff}=t_{lk}^{ff}$
  \end{itemize}
\end{frame}

\begin{frame}{Equilibrium}
\begin{itemize}
  \item Given fundamentals $\{\bar u_i,\bar A_i,t_{kl}^{ff},n_{kl}\}$ with $t_{kl}^{ff}=t_{lk}^{ff}$ on each physical road and parameters $(\theta,\alpha,\beta,\lambda)$, an equilibrium is
  \vspace{-0.5em}\begin{equation}
    \left\{L_i^R,L_i^F,L_{ij},\tau_{ij},t_{kl},\Xi_{kl},\Pi_i,P_j,\chi\right\}
  \vspace{-0.5em}\end{equation}
  such that
  \begin{itemize}
    \item residents choose residence-workplace-route options according to \eqref{eq:pi}
    \item bilateral flows satisfy commuting gravity \eqref{eq:gravity} and market clearing \eqref{eq:clear1}--\eqref{eq:clear2}
    \item amenities and productivity satisfy \eqref{eq:agglom}
    \item bilateral costs are generated by the directed road network \eqref{eq:inverse}
    \item directed link traffic and directed edge times are mutually consistent \eqref{eq:Xi_c}--\eqref{eq:c_Xi}
  \end{itemize}
  \item Existence and uniqueness: see Allen and Arkolakis (2022)
\end{itemize}
\end{frame}

\begin{frame}{Tidal Policy and Observed-Anchored Congestion}
\begin{itemize}
  \item \textbf{Instrument.} On a physical road $\{k,l\}$ hold $N_{kl}=n_{kl}+n_{lk}$ fixed and shift $\Delta$ lanes
  \vspace{-0.4em}\begin{equation}
    n'_{kl}=n_{kl}+\Delta\ (\text{slow}),\qquad n'_{lk}=n_{lk}-\Delta\ (\text{fast}).
  \end{equation}
  \item \textbf{Congestion anchored at the observed allocation}, the empirical form of $t_{kl}=t_{kl}^{ff}(\Xi_{kl}/n_{kl})^{\lambda}$:
  \vspace{-0.4em}\begin{equation}\label{eq:anchor}
    t_{kl}=\max\!\Big( t_{kl}^{\,\mathrm{ff}},\;\;
    t_{kl}^{\,\mathrm{obs}}\big[(\Xi_{kl}/\Xi_{kl}^{\mathrm{obs}})(n_{kl}^{\mathrm{obs}}/n_{kl})\big]^{\lambda}\Big)
  \end{equation}
  \begin{itemize}\footnotesize
    \item Anchored at $t^{\mathrm{obs}}$: baseline \good{matches the data in the median} ($t/t^{\mathrm{obs}}=1.00$); $\sim$8\% of edges hit the $t^{\mathrm{ff}}$ floor where $\tau^{\mathrm{obs}}<\tau^{\mathrm{ff}}$.
    \item $\Xi^{\mathrm{obs}}$ is model traffic at the observed allocation; the policy enters only through capacity $n_{kl}$, with a free-flow floor $t^{\mathrm{ff}}$.
  \end{itemize}
\end{itemize}
\end{frame}

\begin{frame}{Solving the Counterfactual in Levels}
\begin{itemize}
  \item Fundamentals $(\bar u_i,\bar A_j)$ are inverted from the observed city and \good{held fixed} across all scenarios.
  \item A policy changes lane capacity $n_{kl}$; the model recomputes the \good{congested equilibrium} (edge times $t_{kl}$, the sorting of residents and jobs, traffic $\Xi_{kl}$, and market access) and compares it to the baseline solved the same way.
  \item Counterfactual effects are read off as the change between the two equilibria, both solved in levels.
  \item \textbf{Welfare} $\;\bar W_\theta=\sum_i u_i\,\Pi_i^{-\theta}$, with $u_i=\bar u_i\,(L_i^R)^{\beta}$.
\end{itemize}
\end{frame}

\begin{frame}{Short Run vs.\ Long Run}
\begin{itemize}
  \item The same level solver, two horizons:
  \begin{itemize}\footnotesize
    \item \good{Short run (ST)}: pin $L^R,L^F$ at the baseline; only $\tau$ (congestion) re-equilibrates.
    \item \good{Long run (LT)}: full residence / workplace sorting \emph{and} congestion feedback.
  \end{itemize}
  \item \textbf{Decompose} the OD-weighted commuting-cost change ($s$ = OD shares):
  \vspace{-0.4em}\begin{equation}
    \Delta\bar\tau^{LT}=\underbrace{\textstyle\sum s^0(\tau^{ST}\!-\tau^0)}_{\text{ST direct}}
    +\underbrace{\textstyle\sum s^0(\tau^{LT}\!-\tau^{ST})}_{\text{GE wedge}}
    +\underbrace{\textstyle\sum(s^{LT}\!-s^0)\tau^{LT}}_{\text{sorting}}
  \end{equation}
  \item Run the baseline and each policy through the same solver; report ST, GE, and sorting effects on $\bar\tau$ and $\bar W$.
\end{itemize}
\end{frame}

% =====================================================================
\section{Quantification}\label{sec:quant}
% =====================================================================
\begin{frame}{Calibration / Main Specification}
\begin{center}
\begin{tabular}{llll}
\toprule
Parameter & Value & Interpretation & Source \\
\midrule
$\theta$ & $6.83$ & commuting-cost elasticity & AA / Ahlfeldt et al. \\
$\alpha$ & $\good{+0.05}$ & production agglomeration & literature gross spillover \\
$\beta$ & $\bad{-0.10}$ & residential crowding disamenity & AA Seattle \\
$\lambda$ & $1.1$ & congestion (edge time vs traffic per lane) & Barwick; Yang et al. \\
\bottomrule
\end{tabular}
\end{center}
\vspace{0.4em}
\begin{itemize}\small
  \item \textbf{Congestion technology} is the power form $t_{kl}=t_{kl}^{ff}(\Xi_{kl}/n_{kl})^{\lambda}$, $\lambda=1.1$ (not a BPR functional form).
  \item \textbf{Policy.} Top-$100$ reciprocal pairs by the \emph{BPR congestion score}, $\Delta n=1$ (slow $+1$ / fast $-1$).
  \begin{itemize}\footnotesize
    \item Score $=\Xi_s(t_s^{\mathrm{obs}}-t_s^{ff})-\Xi_f(t_f^{\mathrm{obs}}-t_f^{ff})$: baseline flow $\times$ congestion delay, slow minus fast.
    \item Ranks corridors by the first-order short-run relief from moving one lane to the slow side.
  \end{itemize}
  \item \textbf{Horizons.} \good{Short run} pins $L^R,L^F$ at baseline; \good{long run} allows full sorting plus congestion feedback.
\end{itemize}
\end{frame}

\begin{frame}{Baseline Fundamentals (inverted)}
\begin{columns}[T,onlytextwidth]
\column{0.5\textwidth}
\centering \scriptsize Residence amenity $\bar u_i$\\
\includegraphics[width=\linewidth,height=0.52\textheight,keepaspectratio]{fund_ubar.png}
\column{0.5\textwidth}
\centering \scriptsize Workplace productivity $\bar A_j$\\
\includegraphics[width=\linewidth,height=0.52\textheight,keepaspectratio]{fund_abar.png}
\end{columns}
\vspace{0.1em}
\begin{itemize}\footnotesize
  \item $\bar u,\bar A$ are inverted to match observed residence and workplace shares at observed travel times (Assumption: observed allocation is an equilibrium).
  \item $\mathrm{corr}(\bar A_j,L_j^F)=\good{0.67}$ vs.\ $\mathrm{corr}(\bar u_i,L_i^R)=0.30$: jobs pinned by productivity, residence is access-driven.
  \item Treated corridors land on \bad{high-flow employment bottlenecks} (BPR score vs.\ slow-direction $\Xi_s$: corr.\ $\approx+0.91$ on selected top-100 pairs; gate5 pair-selection QA). \app{appfund}
\end{itemize}
\end{frame}

\begin{frame}{Result 1. Tidal Lanes: ST / GE / Sorting Cascade}
\framesubtitle{Average commuter route-cost change (min/person), baseline-OD weighted; main spec ($\alpha{=}{+}.05,\beta{=}{-}.10$).}
\begin{center}\normalsize
\begin{tabular}{lc}
\toprule
Channel & $\Delta$ commute \\
\midrule
ST direct (fixed sorting) & \good{$-1.63$ s} \\
$+$ network GE wedge & $+0.99$ s \\
$+$ residence--workplace sorting & \bad{$+0.71$ s} \\
\midrule
\textbf{Long-run total} & \textbf{$+0.07$ s} \;($+95$ person-min/day) \\
\bottomrule
\end{tabular}
\end{center}
\vspace{0.2em}
\begin{itemize}\small
  \item ST: holding jobs and homes fixed, the corridor \good{really saves time} ($-1.63$ s; slow-direction $\tau$ down \good{$13.6\%$} median).
  \item \textbf{Exposure:} \good{56.4\%} of baseline OD commuter mass lies on shortest paths through treated corridors (path-based, not only endpoints).
  \item LT: \bad{sorting almost exactly offsets it}, so commute nets to $\approx 0$.
  \item Not ``corridor useless''; this is a long-run sorting-direction problem.
\end{itemize}
\end{frame}

\begin{frame}{Result 2. Welfare: ST $\uparrow$, LT $\downarrow$, Loss via $\Pi$}
\begin{columns}[T,onlytextwidth]
\column{0.46\textwidth}
\centering\small
\begin{tabular}{lc}
\toprule
$\bar W$ & change \\
\midrule
Short run & \good{$+0.94\%$} \\
Long run & \bad{$-0.85\%$} \\
\bottomrule
\end{tabular}
\vspace{0.6em}

\scriptsize $\Pi^{-\theta}$ four-layer ($u^0$-weighted, $\Delta\bar W_\theta$):
\begin{tabular}{lc}
\toprule
\textcircled{1} ST $\tau$ @ $\ell^0$ & \good{$+4.60$} \\
\textcircled{2} LT $\tau$ network GE & \bad{$-3.70$} \\
\textcircled{3} $+\Delta\ell^F$ (jobs) & \bad{$-2.90$} \\
\textcircled{4} $+\Delta\ell^R$ (homes) & \bad{$-1.91$} \\
\bottomrule
\end{tabular}
\column{0.52\textwidth}
\centering \scriptsize $\Delta\log\Pi$ (market access), treated endpoints marked\\
\includegraphics[width=\linewidth,height=0.62\textheight,keepaspectratio]{sort_Pi.png}
\end{columns}
\vspace{0.2em}
\begin{itemize}\footnotesize
  \item Loss is \bad{almost entirely through market access $\Pi$} (network GE plus job/home reallocation), not amenity $u$.
\end{itemize}
\end{frame}

\begin{frame}{Result 3. Long-Run Sorting (maps)}
\begin{columns}[T,onlytextwidth]
\column{0.5\textwidth}
\centering \scriptsize $\Delta\log\ell^R$ (residence)\\
\includegraphics[width=\linewidth,height=0.6\textheight,keepaspectratio]{sort_lR.png}
\column{0.5\textwidth}
\centering \scriptsize $\Delta\log\ell^F$ (workplace)\\
\includegraphics[width=\linewidth,height=0.6\textheight,keepaspectratio]{sort_lF.png}
\end{columns}
\vspace{0.1em}
\begin{itemize}\footnotesize
  \item Treated endpoints gain residents and jobs (net change small) but reshuffle the network.
  \item \bad{$\beta<0$} pushes homes off the densifying hubs.
\end{itemize}
\end{frame}

\begin{frame}{Result 4. Why Commute $\neq$ Welfare ($\alpha$--$\beta$ grid)}
\framesubtitle{BPR-score selection, top 100, $\lambda=1.1$. ST direct $\equiv-1.63$ s and $\bar W$ ST $\equiv+0.94\%$ throughout.}
\begin{columns}[T,onlytextwidth]
\column{0.5\textwidth}
\centering Vary $\alpha$ ($\beta=-0.10$)\\[2pt]
\scriptsize
\begin{tabular}{cccc}
\toprule
$\alpha$ & sort (s) & per-min/day & $\bar W$ LT \\
\midrule
$-0.12$ & $-3.63$ & $-6417$ & $-0.32\%$ \\
$0$     & $-1.33$ & $-2960$ & $-0.60\%$ \\
$+0.03$ & $-0.25$ & $-1343$ & $-0.73\%$ \\
$\mathbf{+0.05}$ & $\mathbf{+0.71}$ & $\mathbf{+95}$ & $\mathbf{-0.85\%}$ \\
$+0.10$ & $+4.88$ & $+6340$ & $-1.34\%$ \\
\bottomrule
\end{tabular}
\column{0.5\textwidth}
\centering Vary $\beta$ ($\alpha=+0.05$)\\[2pt]
\scriptsize
\begin{tabular}{cccc}
\toprule
$\beta$ & sort (s) & per-min/day & $\bar W$ LT \\
\midrule
$\mathbf{-0.10}$ & $\mathbf{+0.71}$ & $\mathbf{+95}$ & $\mathbf{-0.85\%}$ \\
$0$     & $-1.57$ & $-3173$ & $-1.11\%$ \\
$+0.05$ & $-4.64$ & $-7579$ & $-1.49\%$ \\
$+0.10$ & \multicolumn{3}{c}{\emph{not converged (appendix only)}} \\
\bottomrule
\end{tabular}
\end{columns}
\vspace{0.4em}
\begin{itemize}\footnotesize
  \item \bad{$\bar W$ LT $<0$ in every converged cell}, monotone in both $\alpha$ and $\beta$.
  \item Commute-LT sign is the \good{$\alpha$--$\beta$ knife-edge}: at $\alpha{=}{-}.12$ commute improves a lot ($-6417$) yet $\bar W$ still falls.
  \item \textbf{Time savings $\neq$ welfare.} \app{appgrid}\;\app{appgridext}
\end{itemize}
\end{frame}

\begin{frame}{Result 5. Exogenous vs.\ Endogenous Feedback}
\begin{center}\small
\begin{tabular}{lcc}
\toprule
Counterfactual & exogenous $\tau$ & endogenous \\
\midrule
CF1 symmetric cost & $-2.77\%$ & $\approx 0\%$ \\
CF3 uniform $\times1.10$ & $-0.40\%$ & $-0.12\%$ \\
Tidal top100 (LT) & n/a & $\bad{-0.85\%}$ \\
\bottomrule
\end{tabular}
\end{center}
\vspace{0.4em}
\begin{itemize}\small
  \item Anchored congestion \good{re-floods} pure-$\tau$ shocks: symmetrization and uniform speed-ups are largely undone as $\Xi$ re-equilibrates.
  \item \good{Tidal works because it changes capacity $n$} (which enters the anchor), not a $\tau$ multiplier. That is the theoretical reason the instrument bites. \app{appexog}
\end{itemize}
\end{frame}

\begin{frame}[label=result6]{Result 6. Policy Comparison (ST / LT $\bar W$)}
\framesubtitle{Same gate5 mainline; endogenous anchored congestion ($\lambda=1.1$). Lane-km matched to P2/P3 expansion $\approx 1{,}092$ km.}
\begin{center}\scriptsize
\begin{tabular}{lcccr}
\toprule
Policy & ST $\bar W$ & LT $\bar W$ & lane-km & Role \\
\midrule
P6 first-best pricing & \good{$+15.3\%$} & $-0.49\%$ & 0 & welfare upper bound \\
P3 expansion (qdelay top100) & \good{$+1.65\%$} & \bad{$-1.39\%$} & 1{,}092 & large ST, LT sorting cost \\
P2 expansion (same corridor) & $+0.99\%$ & $-0.92\%$ & 1{,}092 & matched lane budget \\
\textbf{P1 tidal (top100)} & \good{$+0.94\%$} & \bad{$-0.85\%$} & \good{\textbf{0}} & \textbf{zero net capacity} \\
P7 cordon (q0.9) & $-0.19\%$ & $+0.09\%$ & 0 & weak LT gain \\
P4 uniform (matched km) & $-0.02\%$ & $-0.02\%$ & 1{,}092 & undirected benchmark \\
\bottomrule
\end{tabular}
\end{center}
\vspace{0.3em}
\begin{itemize}\footnotesize
  \item \textbf{Takeaway:} pricing dominates on $\bar W$; undirected expansion buys lane-km but LT welfare stays negative once sorting feeds back.
  \item \textbf{Tidal niche:} \good{zero lane-km}, \good{directional} capacity shift; ST welfare near expansion, without building roads. \app{apppolicy}
\end{itemize}
\end{frame}

\begin{frame}{Conclusion}
\begin{enumerate}
  \item \textbf{Directional congestion is a widespread, first-order urban feature.}
  \begin{itemize}\footnotesize
    \item Across top-50 cities, \good{$\sim$52\%} of weighted OD pairs are $\geq10\%$ asymmetric (median ratio $0.90$); flow-weighted Shenzhen is stronger ($0.82$); it persists down to Beijing road segments.
  \end{itemize}
  \item \textbf{Ignoring direction mis-measures welfare.}
  \begin{itemize}\footnotesize
    \item A symmetric commuting cost smooths the bottlenecks and \bad{overstates} welfare (CF1 symmetric $-2.77\%$ under exogenous $\tau$ vs $\approx0\%$ once congestion re-equilibrates), and mis-allocates real commute exposure.
    \item Capturing the policy needs \emph{both} directional cost and endogenous feedback.
  \end{itemize}
  \item \textbf{Tidal lanes: cost-effective and directional, not welfare-maximizing.}
  \begin{itemize}\footnotesize
    \item Short run they save time ($-1.63$ s, $\bar W$ $+0.94\%$); long run, sorting capitalizes the gain into market access, so commute nets to $\approx0$ ($+95$/day) and $\bar W$ turns \bad{$-0.85\%$} (loss via $\Pi$, robust across $\alpha,\beta$).
    \item Versus pricing and matched expansion (Result~6): \good{zero net lane-km} and directional targeting; not the aggregate-welfare maximum.
  \end{itemize}
\end{enumerate}
\end{frame}

% =====================================================================
\appendix
\section{Appendix}
% =====================================================================
\begin{frame}[label=appdata]{Appendix: Data Construction}
\begin{itemize}\small
  \item \textbf{Links.}
  \begin{itemize}\footnotesize
    \item Raw road geometry to links between adjacent nodes; drop duplicate/unusable records.
    \item Each physical road kept as symmetric infrastructure with two directed flows.
  \end{itemize}
  \item \textbf{Speeds.}
  \begin{itemize}\footnotesize
    \item $v_{klh}=\mathrm{dist}_{kl}/\mathrm{time}_{klh}$, inverse speed $m_{klh}=v_{klh}^{-1}$; AM peak selected.
  \end{itemize}
  \item \textbf{Grid and OD.}
  \begin{itemize}\footnotesize
    \item Locations aggregated to square cells; OD commuting linked to the network.
    \item Restrict to the largest connected component, with a detour gate to drop implausible links.
  \end{itemize}
  \item \textbf{Lanes.}
  \begin{itemize}\footnotesize
    \item $n_{kl}$ from a corrected road-segment lane lookup; $1074$ unique values, $p50\approx3.4$.
  \end{itemize}
\end{itemize}
\centering\back{data}
\end{frame}

\begin{frame}[label=appfund]{Appendix: Inverting Fundamentals}
\begin{itemize}\small
  \item Treat the observed allocation as an equilibrium; recover $(\bar u_i^\theta,\bar A_j^\theta)$ by alternating updates with geometric-mean normalization over identifying (positive-on-both-sides) nodes.
  \item $\kappa$ rescaling enforces the spectral-radius condition $\rho(A)<1$; counterfactual ratios are invariant to $\kappa$ (held fixed across baseline / anchor / CF).
  \item Fit: $\mathrm{corr}(\bar A_j,L_j^F)=0.67$, $\mathrm{corr}(\bar u_i,L_i^R)=0.30$; welfare geographically diffuse (top-5 nodes $=4.2\%$ of $\bar W$ weight).
\end{itemize}
\centering\back{data}
\end{frame}

\begin{frame}[label=appestim]{Appendix: Estimation / Solver}
\begin{itemize}\small
  \item \textbf{Calibration (external):} $(\theta,\alpha,\beta)$ from AA / literature; $\lambda_{\mathrm{phys}}=1.1$ from Barwick et al.\ (2024) mapping---\emph{not} structurally estimated on Beijing flows (no directed traffic panel yet).
  \item Anchored implementation: $t_{kl}=\max\big(t_{kl}^{ff},\,t_{kl}^{\mathrm{obs}}[(\Xi_{kl}/\Xi_{kl}^{\mathrm{obs}})(n_{kl}^{\mathrm{obs}}/n_{kl})]^{\lambda}\big)$; median $t/t^{\mathrm{obs}}\approx 1.00$ (see free-flow floor note above).
  \item Level-form solver: outer loop on $t_{kl}$, inner alternating solve of $(L^R,L^F)$ then market access $(\Pi,P)$; damped log-update.
  \item Stability needs $\alpha\theta<1$ and $|\alpha\beta\theta^2|<1$ (satisfied on the grid); $\alpha>0$ uses $\mathrm{damping}=0.10$, $\mathrm{max\_outer}=500$, $\mathrm{tol}=10^{-4}$.
  \item Full parameter map: \texttt{Documents/L2\_notes/TidalLanes\_estimation\_notes.pdf}.
\end{itemize}
\centering\back{data}
\end{frame}

\begin{frame}[label=appgrid]{Appendix: Full $\alpha$--$\beta$ Sensitivity (9 cells)}
\scriptsize
\begin{center}
\begin{tabular}{ccccccc}
\toprule
$\alpha$ & $\beta$ & ST dir (s) & sort (s) & LT (s) & per-min/day & $\bar W$ LT \\
\midrule
$-0.12$ & $-0.10$ & $-1.63$ & $-3.63$ & $-4.37$ & $-6417$ & $-0.32\%$ \\
$0.00$ & $-0.10$ & $-1.63$ & $-1.33$ & $-2.02$ & $-2960$ & $-0.60\%$ \\
$0.02$ & $-0.10$ & $-1.63$ & $-0.65$ & $-1.32$ & $-1942$ & $-0.69\%$ \\
$0.03$ & $-0.10$ & $-1.63$ & $-0.25$ & $-0.91$ & $-1343$ & $-0.73\%$ \\
$\mathbf{0.05}$ & $\mathbf{-0.10}$ & $\mathbf{-1.63}$ & $\mathbf{+0.71}$ & $\mathbf{+0.06}$ & $\mathbf{+95}$ & $\mathbf{-0.85\%}$ \\
$0.10$ & $-0.10$ & $-1.63$ & $+4.88$ & $+4.32$ & $+6340$ & $-1.34\%$ \\
$0.05$ & $0.00$ & $-1.63$ & $-1.57$ & $-2.16$ & $-3173$ & $-1.11\%$ \\
$0.05$ & $0.05$ & $-1.63$ & $-4.64$ & $-5.16$ & $-7579$ & $-1.49\%$ \\
$0.05$ & $0.10$ & \multicolumn{5}{c}{\emph{not converged} (median $\tau/\tau^{\mathrm{obs}}\approx1.35$); diagnostic only} \\
\bottomrule
\end{tabular}
\end{center}
\centering\back{data}
\end{frame}

\begin{frame}[label=appgridext]{Appendix: Extended $\alpha$--$\beta$ Grid (+9 cells)}
\scriptsize
\begin{itemize}\footnotesize
  \item Adds targeted cells (negative $\alpha$, $\beta\in\{-0.12,-0.08,-0.05\}$, etc.) to search for $\bar W_{\mathrm{LT}}>0$; all converged cells remain \bad{negative}.
  \item Best LT welfare in extended grid: $(\alpha,\beta)=(-0.12,-0.12)$ $\Rightarrow$ $\bar W_{\mathrm{LT}}=\bad{-0.32\%}$ (person-min/day $-6{,}197$).
\end{itemize}
\vspace{0.2em}
\begin{center}
\begin{tabular}{lcccc}
\toprule
$(\alpha,\beta)$ & sort (s) & per-min/day & $\bar W$ LT & conv. \\
\midrule
$(-0.12,-0.12)$ & $-3.47$ & $-6{,}197$ & \textbf{$-0.32\%$} & yes \\
$(-0.12,-0.10)$ & $-3.63$ & $-6{,}417$ & $-0.32\%$ & yes \\
$(+0.05,-0.10)$ & $+0.71$ & $+95$ & $-0.85\%$ & yes (main) \\
$(+0.05,+0.10)$ & --- & --- & --- & \emph{no} \\
\bottomrule
\end{tabular}
\end{center}
\centering\back{data}
\end{frame}

\begin{frame}[label=apppolicy]{Appendix: Policy Comparison (P1--P7)}
\scriptsize
\begin{center}
\begin{tabular}{lccccc}
\toprule
Policy & horizon & $\bar W$ chg & res.\ diss.\ (\%) & wp.\ diss.\ (\%) & lane-km \\
\midrule
P1 tidal top100 & ST & $+0.94\%$ & $\approx 0$ & $\approx 0$ & 0 \\
P1 tidal top100 & LT & $-0.85\%$ & 0.86 & 1.79 & 0 \\
P6 first-best pricing & ST & $+15.3\%$ & $\approx 0$ & $\approx 0$ & 0 \\
P6 first-best pricing & LT & $-0.49\%$ & 5.87 & 10.33 & 0 \\
P3 qdelay expansion & ST & $+1.65\%$ & $\approx 0$ & $\approx 0$ & 1{,}092 \\
P3 qdelay expansion & LT & $-1.39\%$ & 1.10 & 1.92 & 1{,}092 \\
P2 same-corridor exp. & LT & $-0.92\%$ & 0.94 & 1.70 & 1{,}092 \\
\bottomrule
\end{tabular}
\end{center}
\vspace{0.2em}
\begin{itemize}\footnotesize
  \item Source: \texttt{policy\_comparison\_gate5\_bpr\_pos05\_v1/results\_by\_policy.csv}.
  \item Mean ST/LT $\bar W$ ranking: pricing $\gg$ qdelay expansion $>$ tidal $\approx$ same-corridor expansion.
\end{itemize}
\centering\back{result6}
\end{frame}

\begin{frame}[label=appexog]{Appendix: Exogenous vs.\ Endogenous (full)}
\scriptsize
\begin{center}
\begin{tabular}{llcc}
\toprule
Family & CF & exog $\tau$ ($\bar W$) & endog ($\bar W$) \\
\midrule
structural & CF1 symmetric & $-2.77\%$ & $\approx 0\%$ \\
structural & CF3 uniform $\times1.10$ & $-0.40\%$ & $-0.12\%$ \\
structural & CF3 uniform $\times1.25$ & $-1.38\%$ & $-0.43\%$ \\
place-based & uniform speed $+10\%$ & $-0.46\%$ & $+0.03\%$ \\
place-based & hub access $+15\%$ & $-2.68\%$ & $+0.01\%$ \\
place-based & hub productivity $+10\%$ & $-1.45\%$ & $-1.27\%$ \\
place-based & outer-ring amenity $+10\%$ & $+0.57\%$ & $+0.49\%$ \\
\bottomrule
\end{tabular}
\end{center}
\vspace{0.3em}
\begin{itemize}\footnotesize
  \item Transport ($\tau$) shocks damp out under endogenous congestion; fundamental (productivity/amenity) shocks pass through.
\end{itemize}
\centering\back{data}
\end{frame}

\begin{frame}{Appendix: Robustness Notes}
\begin{itemize}\small
  \item \textbf{Selection.}
  \begin{itemize}\footnotesize
    \item BPR congestion score is main; the asymmetry-ratio rule ($\tau_{\text{slow}}/\tau_{\text{fast}}$) is robustness (overlap $\sim$38/100).
    \item Under the asymmetry-ratio rule, the main spec gives ST $-0.79$ s, $\bar W$ LT $-0.37\%$ (same signs, smaller magnitudes).
  \end{itemize}
  \item \textbf{Congestion band.}
  \begin{itemize}\footnotesize
    \item Main $\lambda=1.1$; robustness $\lambda=0.57$ (legacy AA slope band): top100 LT person-min $\approx -1{,}168$/day vs.\ $+95$ at $\lambda=1.1$.
  \end{itemize}
  \item \textbf{Convergence.}
  \begin{itemize}\footnotesize
    \item All main cells converge to $\mathrm{tol}=10^{-4}$; residual dominated by $\sim$8\% of edges with $\tau^{\mathrm{obs}}<\tau^{\mathrm{ff}}$ (free-flow floor), independent of policy.
    \item $\beta\geq+0.10$ fails to converge under $\lambda=1.1$.
  \end{itemize}
\end{itemize}
\centering\back{data}
\end{frame}

\end{document}
